\documentclass[11pt,a4paper]{article}

% --- Packages ---
\usepackage[utf8]{inputenc}
\usepackage[T1]{fontenc}
\usepackage[margin=2.5cm]{geometry}
\usepackage{amsmath,amssymb,amsfonts}
\usepackage{mathtools}
\usepackage{empheq}
\usepackage{siunitx}
% siunitx v3 removed \SI and \SIrange; v2 provided them. Add a compatibility
% shim so the manual compiles under both versions.
\providecommand{\SI}[2]{\qty{#1}{#2}}
\providecommand{\SIrange}[3]{\qtyrange{#1}{#2}{#3}}
\usepackage{graphicx}
\usepackage{xcolor}
\usepackage{tcolorbox}
\tcbuselibrary{breakable,skins}
\usepackage{booktabs}
\usepackage{array}
\usepackage{enumitem}
\usepackage{hyperref}
\usepackage{titlesec}
\usepackage{fancyhdr}
\usepackage{tikz}
\usetikzlibrary{arrows.meta,positioning,calc,decorations.pathreplacing,patterns}

% --- Colours ---
\definecolor{deepblue}{RGB}{31,56,100}
\definecolor{midblue}{RGB}{46,117,182}
\definecolor{accentpink}{RGB}{255,107,157}
\definecolor{accentcyan}{RGB}{0,180,210}
\definecolor{warnorange}{RGB}{255,140,30}
\definecolor{dangerred}{RGB}{200,50,70}
\definecolor{boxgrey}{RGB}{245,247,250}
\definecolor{rulegrey}{RGB}{210,215,225}

% --- Hyperref setup ---
\hypersetup{
  colorlinks=true,
  linkcolor=deepblue,
  citecolor=midblue,
  urlcolor=midblue,
  pdftitle={Virtual Tri-HB Integrated -- Student Manual},
  pdfauthor={Monash Tri-HB Group}
}

% --- Title formatting ---
\titleformat{\section}
  {\Large\bfseries\color{deepblue}}
  {\thesection.}{0.6em}{}
  [\vspace{2pt}{\color{rulegrey}\hrule height 0.6pt}]
\titleformat{\subsection}
  {\large\bfseries\color{midblue}}
  {\thesubsection}{0.6em}{}
\titleformat{\subsubsection}
  {\normalsize\bfseries\color{midblue}}
  {\thesubsubsection}{0.6em}{}

% --- Headers and footers ---
\pagestyle{fancy}
\fancyhf{}
\fancyhead[L]{\small\color{deepblue}Virtual Tri-HB Integrated -- Student Manual}
\fancyhead[R]{\small\color{deepblue}\nouppercase{\rightmark}}
\fancyfoot[C]{\small\thepage}
\renewcommand{\headrulewidth}{0.4pt}
\renewcommand{\headrule}{\color{rulegrey}\hrule height 0.4pt}

% --- Custom boxes ---
\newtcolorbox{conceptbox}[1][]{
  colback=boxgrey, colframe=midblue,
  boxrule=0.6pt, arc=2pt,
  fonttitle=\bfseries\color{white},
  coltitle=white, colbacktitle=midblue,
  title={#1}, breakable,
  left=8pt, right=8pt, top=6pt, bottom=6pt
}
\newtcolorbox{stepbox}[1][]{
  colback=white, colframe=accentcyan,
  boxrule=1pt, arc=1pt,
  fonttitle=\bfseries\color{white},
  coltitle=white, colbacktitle=accentcyan,
  title={#1}, breakable,
  left=8pt, right=8pt, top=6pt, bottom=6pt
}
\newtcolorbox{warningbox}[1][]{
  colback=boxgrey, colframe=warnorange,
  boxrule=1pt, arc=1pt,
  fonttitle=\bfseries\color{white},
  coltitle=white, colbacktitle=warnorange,
  title={#1}, breakable,
  left=8pt, right=8pt, top=6pt, bottom=6pt
}
\newtcolorbox{newfeature}[1][]{
  colback=boxgrey, colframe=accentpink,
  boxrule=1pt, arc=1pt,
  fonttitle=\bfseries\color{white},
  coltitle=white, colbacktitle=accentpink,
  title={#1}, breakable,
  left=8pt, right=8pt, top=6pt, bottom=6pt
}

% --- Mathematical shortcuts ---
\newcommand{\sig}{\sigma}
\newcommand{\eps}{\varepsilon}
\newcommand{\epsdot}{\dot\eps}
\newcommand{\Cz}{C_0}
\newcommand{\Ab}{A_b}
\newcommand{\As}{A_s}
\newcommand{\Ls}{L_s}
\newcommand{\sigprop}{\sig_{\text{prop}}}
\newcommand{\sigpeak}{\sig_{\text{peak}}}

% --- Title page content ---
\title{
  \vspace{-1.5cm}
  {\Huge\bfseries\color{deepblue} Virtual Tri-HB}\\[6pt]
  {\Large\bfseries\color{midblue} Integrated Workspace -- Student User Manual}\\[14pt]
  {\large\itshape Four Loading Modes plus Experimental Analysis,\\
  Wave, Stress-Path, Damage and DEM Validation}
}
\author{
  Monash University $\cdot$ Tri-HB Research Group
}
\date{Companion to \texttt{tri\_hb\_integrated.py} (Integrated Edition)}

\begin{document}

% =====================================================
% TITLE PAGE
% =====================================================
\maketitle
\thispagestyle{empty}

\vspace{1.0cm}
\begin{center}
\begin{tcolorbox}[colback=boxgrey, colframe=midblue, width=0.88\textwidth, arc=3pt]
\begin{center}
{\large\bfseries\color{deepblue} About This Manual}\\[6pt]
\end{center}
This manual is written for students learning dynamic mechanics testing through the
Virtual Tri-HB simulator. It explains the physics, equations, and analysis
procedures for each of the four loading modes available in the simulator, and for
the experimental data reducer, wave model, stress-path interpreter, and damage
model that surround it. Every equation is derived from first principles, with the
steps a student would actually need to follow in lab work or in homework problems.

The application is now delivered as a single \textbf{integrated Streamlit
workspace} (\texttt{tri\_hb\_integrated.py}) that connects four analysis steps
under one shared experimental setup:

\begin{enumerate}[leftmargin=2.4em, itemsep=2pt]
  \item \textbf{Step 1 -- Setup, simulator and data:} configure the bar, specimen,
        material, and loading; run the simulator; \emph{or} reduce uploaded bar
        gauge / direct stress-strain data.
  \item \textbf{Step 2 -- Wave model:} pulse timing, travel time, equilibrium
        window, and the wave-superposition / damage-memory regime classifier.
  \item \textbf{Step 3 -- Stress path and analysis:} $p$--$q$--$\theta$ stress
        paths, invariant histories, and comparison against reduced test data.
  \item \textbf{Step 4 -- Damage model and DEM validation:} damage evolution,
        stiffness and wave-speed degradation, energy balance, and synthetic DEM
        descriptors.
\end{enumerate}

The four loading modes available in Step 1 are unchanged from earlier editions:
\begin{enumerate}[leftmargin=2.4em, itemsep=2pt]
  \item Gas-Gun uniaxial impact (classical SHPB)
  \item Electromagnetic uniaxial impact (clean half-sine pulse)
  \item Electromagnetic asynchronous triaxial impact (stress-path control)
  \item Electromagnetic symmetric multidirectional impact (wave superposition)
\end{enumerate}

All four modes accept independent static pre-stress on \emph{three} orthogonal
axes ($\sig_1$, $\sig_2$, $\sig_3$), applied through the hydraulic cylinders
before the dynamic pulse arrives.
\end{tcolorbox}
\end{center}

\newpage
\tableofcontents
\newpage

% =====================================================
% PART 1 -- FOUNDATIONS
% =====================================================
\section{Foundations: What Every Tri-HB Test Has in Common}

Before separating the four loading modes, students should understand the building
blocks shared by all of them. Every Tri-HB test, regardless of mode, involves the
following physical elements: a long elastic bar (or several), a stress wave
travelling along that bar, a small specimen at one or more bar interfaces, and
strain gauges that record the wave passage. The mathematics of the bar and the
mathematics of the specimen are kept separate, and the analysis bridges them
through force and velocity continuity at the interfaces.

\subsection{The 1-D Wave Equation in the Bars}

Each pressure bar in the Tri-HB system is a long, slender rod of a high-strength
steel (typically 42CrMo). Provided the wavelength of the propagating pulse
is much greater than the bar diameter, the axial displacement $u(x,t)$ obeys the
one-dimensional elastic wave equation:
\begin{equation}
  \frac{\partial^2 u}{\partial t^2} = \Cz^2 \frac{\partial^2 u}{\partial x^2},
  \qquad \Cz = \sqrt{\frac{E_b}{\rho_b}}.
  \label{eq:wave}
\end{equation}
For 42CrMo steel ($E_b = \SI{210}{GPa}$, $\rho_b = \SI{7850}{kg/m^3}$),
the bar wave speed is $\Cz \approx \SI{5172}{m/s}$. The Step~1 sidebar exposes
$E_b$ and $\Cz$ as editable number inputs so a different bar material or a
dispersion-corrected calibrated $\Cz$ can be entered without code changes.

The general solution of \eqref{eq:wave} is the d'Alembert form:
\begin{equation}
  u(x,t) = f(x - \Cz t) + g(x + \Cz t),
\end{equation}
where $f$ is a wave travelling in the $+x$ direction and $g$ is a wave travelling
in the $-x$ direction. From this, the axial strain and particle velocity are:
\begin{equation}
  \eps(x,t) = \frac{\partial u}{\partial x},
  \qquad
  v(x,t) = \frac{\partial u}{\partial t} = -\Cz\,\eps_+(x,t) + \Cz\,\eps_-(x,t),
\end{equation}
where the subscripts denote the direction of travel. The cornerstone identity of
Hopkinson bar analysis is $|v| = \Cz\,|\eps|$.

\subsection{Force and Velocity at a Specimen Face}

When a wave reaches a specimen-bar interface, two boundary conditions must hold:
\textbf{force continuity} (the bar pushes the specimen with the same force the
specimen pushes back) and \textbf{velocity continuity} (the bar end moves with the
specimen face). For an incident bar carrying an incoming wave $\eps_I$ and a
reflected wave $\eps_R$:
\begin{align}
  F_1(t) &= E_b \Ab \big[\eps_I(t) + \eps_R(t)\big],
  \label{eq:F1}\\
  v_1(t) &= \Cz \big[\eps_I(t) - \eps_R(t)\big].
  \label{eq:v1}
\end{align}
On the opposite specimen face, a transmission bar carries only an outgoing wave
$\eps_T$:
\begin{align}
  F_2(t) &= E_b \Ab\, \eps_T(t),
  \label{eq:F2}\\
  v_2(t) &= \Cz\, \eps_T(t).
  \label{eq:v2}
\end{align}
\textit{Sign convention used throughout the simulator:} $\eps_I > 0$ corresponds
to a compressive incident wave; $\eps_R$ is positive when the reflected wave is
also compressive (which occurs only if the specimen is stiffer than the bar, a
rare case); for normal brittle specimens $\eps_R$ is tensile (negative). The
transmitted wave $\eps_T$ is positive for a compressive specimen response.

\subsection{The Specimen as a Filter Between Two Bar Signals}

The specimen, of length $\Ls$ and cross-section $\As$, is small compared to the
bar wavelength. After several internal reverberations (typically a few
microseconds), the specimen reaches dynamic stress equilibrium where
$F_1 \approx F_2$. Under this condition, the specimen stress and strain rate are:
\begin{empheq}[box=\fbox]{align}
  \sig_s(t) &= \frac{F_1 + F_2}{2 \As}
            = \frac{E_b \Ab}{2 \As}\big[\eps_I + \eps_R + \eps_T\big],
  \label{eq:sigma3wave}\\
  \epsdot_s(t) &= \frac{v_1 - v_2}{\Ls}
              = \frac{\Cz}{\Ls}\big[\eps_I - \eps_R - \eps_T\big],
  \label{eq:epsdot3wave}\\
  \eps_s(t) &= \int_0^t \epsdot_s(\tau)\,d\tau.
  \label{eq:eps3wave}
\end{empheq}
These three equations are the \textbf{three-wave method}. They are universal --
the same equations hold whether the pulse comes from a striker, an
electromagnetic generator, or anything else, and whether the bars are round or
square.

Under perfect equilibrium $\eps_I + \eps_R = \eps_T$, the strain-rate expression
\eqref{eq:epsdot3wave} simplifies to the classical \emph{one-wave} form
$\epsdot_s = -2\Cz\,\eps_R/\Ls$; equivalently the stress collapses to the
transmitted-wave form $\sig_s = (E_b \Ab / \As)\,\eps_T$. The Step~1 reducer
exposes both forms (``Three-wave average'' and ``Transmitted-wave only'') so
students can compare them on the same dataset.

\begin{conceptbox}[Why three-wave analysis works]
The three-wave equations give us the specimen response without ever modelling the
specimen explicitly. They follow purely from continuity of force and velocity at
the interfaces and from how elastic waves propagate in the bars. A student should
remember: \textbf{the bars are diagnostic instruments, not part of the specimen
under test}. Their job is to record what stress and what strain rate the
specimen experienced, nothing more.
\end{conceptbox}

\subsection{The Bar Plasticity Limit}

The three-wave equations assume the bars remain linearly elastic. 42CrMo
steel has a proportional limit $\sigprop \approx \SI{930}{MPa}$. Whenever the
stress carried by any bar exceeds this value, the bar yields plastically, the
wave equation~\eqref{eq:wave} no longer holds, and the analysis becomes invalid.
This single constraint determines the operating envelope of the entire facility.
\textbf{Always check it before running a test.}

\begin{warningbox}[Bar plasticity: the most important diagnostic]
The maximum stress on any bar must satisfy:
\[
  \boxed{\;\sig_{\text{bar}}^{\max} \leq \sigprop = \SI{930}{MPa}\;}
\]
If exceeded, the test is invalid (and the bar may be damaged). The simulator
flags this automatically with three tiers: caution at $0.85\,\sigprop$, warning
at $\sigprop$, critical at the yield stress \SI{1080}{MPa}.
\end{warningbox}

\subsection{Static Triaxial Pre-stress: $\sigma_1$, $\sigma_2$, $\sigma_3$}
\label{subsec:prestress}

Every Tri-HB test allows independent application of static pre-stress on the
three principal axes via the hydraulic cylinders, \emph{before} the dynamic
pulse arrives. This is what distinguishes the Tri-HB from a classical
uniaxial SHPB. In the Step~1 sidebar, three sliders set the static pre-stress:
$\sig_1$ on the X axis (the loading direction), $\sig_2$ on Y, and $\sig_3$ on
Z. The default slider range is \SIrange{0}{50}{MPa} per axis; this is the value
calibrated against the hydraulic loading frame currently in the digital twin
and is the same upper bound for all four modes.

The static and dynamic contributions superpose linearly to give the total
stress state during the event:
\begin{empheq}[box=\fbox]{align}
  \sigma_X^{\text{total}}(t) &= \sigma_1^{\text{static}} + \sigma_X^{\text{dyn}}(t), \\
  \sigma_Y^{\text{total}}(t) &= \sigma_2^{\text{static}} + \sigma_Y^{\text{dyn}}(t), \\
  \sigma_Z^{\text{total}}(t) &= \sigma_3^{\text{static}} + \sigma_Z^{\text{dyn}}(t).
\end{empheq}

The two roles played by pre-stress are subtly different and worth distinguishing:

\begin{itemize}[leftmargin=*, itemsep=4pt]
  \item \textbf{$\sigma_1$ is an axial baseline.} Because X is the loading
        direction, $\sigma_1$ does not strengthen the rock against the dynamic
        pulse -- it shifts the entire $\sigma$--$\eps$ curve upward by a constant
        offset. Increasing $\sigma_1$ adds to the total recorded axial stress
        but does not change the peak strength of the rock.

  \item \textbf{$\sigma_2$ and $\sigma_3$ are confining stresses.} Because Y
        and Z are perpendicular to the X loading, they confine the rock and
        \emph{strengthen} it via the Mohr--Coulomb mechanism:
    \begin{equation}
      \sigma_{\text{peak}}^{\text{rock}}
        = \big(\sigma_{c0} + k_{\text{conf}} \cdot \max(\sigma_2, \sigma_3)\big) \cdot \text{DIF}.
    \end{equation}
        Doubling $\sigma_2$ raises the peak strength of the specimen; doubling
        $\sigma_1$ does not.
\end{itemize}

\begin{conceptbox}[Why the bar gauge does NOT see the pre-stress]
A common student question: if the static pre-stress is held by the hydraulic
cylinders pressing on the bar ends, why don't the strain gauges on the bars
record it? The answer is that the static pre-stress is reacted by the cylinder
piston, the bar, and the specimen \emph{as a single static system at zero
frequency}. The strain gauges are calibrated to the pre-event baseline (which
is subtracted before analysis), so they record only the \emph{change} from
that baseline -- the dynamic component. The wave-analysis equations in this
manual operate exclusively on the dynamic component. The static pre-stress is
added back at the reporting stage when computing the total stress state.
\end{conceptbox}

% =====================================================
% MODE 1 -- GAS-GUN
% =====================================================
\newpage
\section{Mode 1: Gas-Gun Uniaxial Impact (Classical SHPB)}

This is the original mode of operation, and the easiest to understand. A
projectile (the striker) is launched by compressed gas, strikes the end of an
incident bar, and creates a stress pulse that travels into the specimen.

\subsection{Hardware Configuration}

\begin{itemize}[leftmargin=*, itemsep=2pt]
  \item \textbf{Striker bar}: 42CrMo steel, $L = \SI{0.5}{m}$,
        impact velocity up to \SI{50}{m/s}.
  \item \textbf{Incident and transmission bars}: the simulator's
        \emph{Experimental setup} sidebar tab lets the student choose
        \texttt{Square 50 x 50 mm} ($\Ab = \SI{2500}{mm^2}$) or
        \texttt{Round} with user-supplied diameter
        ($\Ab = \pi (d/2)^2$). The classical gas-gun bench uses a round
        $\Phi 40$ bar ($\Ab \approx \SI{1257}{mm^2}$); the digital
        twin's default is the square cross-section because it is shared with
        the EM bench.
  \item \textbf{Absorption bar} and \textbf{moment-trap device}: prevent
        end-reflections and tensile contamination.
  \item \textbf{Static pre-stress}: hydraulic cylinders independently apply
        $\sig_1$ on the X axis, $\sig_2$ on Y, and $\sig_3$ on Z faces of the
        cubic specimen (slider range \SIrange{0}{50}{MPa} per axis in the
        current simulator build). See Section~\ref{subsec:prestress} for the
        role each plays.
\end{itemize}

\subsection{Step 1: Generating the Incident Pulse}

When the striker, moving at velocity $V$, impacts the stationary incident bar,
both bar ends suddenly acquire a common velocity $V/2$ (by momentum and
continuity). The compressive stress generated is given by the impact relation:
\begin{equation}
  \boxed{\;\sigpeak = \tfrac{1}{2}\,\rho_b\,\Cz\,V \;=\; \tfrac{1}{2}\,(E_b/\Cz)\,V\;}
  \label{eq:gungauge}
\end{equation}
The two forms are identical because $E_b = \rho_b\,\Cz^2$; the simulator
internally uses $\tfrac{1}{2}\,(E_b/\Cz)\,V$ so that the editable bar parameters
$E_b$ and $\Cz$ propagate directly into the peak stress.

\begin{stepbox}[Worked numbers: striker velocity $V = \SI{20}{m/s}$]
  \begin{align*}
  \sigpeak &= \tfrac{1}{2}(\SI{210e9}{Pa} / \SI{5172}{m/s})(20\;\mathrm{m/s})
           \\
           &\approx \SI{406}{MPa}.
  \end{align*}
  Pulse duration is set by the round-trip time within the striker:
  \begin{align*}
  \tau &= \frac{2 L_{\text{striker}}}{\Cz} = \frac{2 \times 0.5}{5172}
       \approx \SI{193}{\micro s}.
  \end{align*}
  So a \SI{20}{m/s} impact produces a roughly rectangular pulse of \SI{406}{MPa}
  amplitude lasting \SI{193}{\micro s}.
\end{stepbox}

In practice, a thin disc of soft metal (the \emph{pulse shaper}, typically copper
or rubber, \SIrange{0.5}{2}{mm} thick) is placed at the impact face. Its plastic
deformation smooths the leading edge of the pulse from a near-instantaneous step
to a finite rise of \SIrange{20}{40}{\micro s}, which (a) reduces high-frequency
oscillations from Pochhammer--Chree dispersion, and (b) gives the specimen time
to reach stress equilibrium before failure. The simulator's gas-gun pulse
generator approximates this with a trapezoidal envelope of \SI{20}{\micro s}
rise/fall and a flat plateau between.

\subsection{Step 2: Pulse Propagation and Specimen Response}

The pulse travels along the incident bar at $\Cz$. When it reaches the specimen,
three things happen simultaneously:
\begin{enumerate}[leftmargin=*]
  \item Part of the pulse is \textbf{reflected} back into the incident bar
        ($\eps_R$, normally tensile for brittle specimens).
  \item Part is \textbf{transmitted} through the specimen into the
        transmission bar ($\eps_T$, compressive).
  \item The specimen \textbf{deforms} at strain rate
        $\epsdot_s = (\Cz/\Ls)(\eps_I - \eps_R - \eps_T)$.
\end{enumerate}
The Y and Z output bars carry only the lateral expansion of the specimen
(Poisson plus damage dilation), since no active dynamic pulse is applied to them
in gas-gun mode.

\subsection{Step 3: Reading the Strain Gauge Records}

Strain gauges on the incident bar record both $\eps_I$ (passing on the way in)
and $\eps_R$ (passing on the way back). A second pair of gauges on the
transmission bar records $\eps_T$. To separate the incident from the reflected
signal in time, the gauge is placed at a distance $L_g$ from the specimen such
that:
\begin{equation}
  \frac{L_g}{\Cz} > \tau / 2,
\end{equation}
ensuring no overlap. For $\tau = \SI{193}{\micro s}$ and a typical
$L_g = \SI{1.0}{m}$, the round-trip time is $\SI{387}{\micro s}$, giving safe
separation.

\subsection{Step 4: Wave Analysis (the Three-Wave Method)}

Apply equations~\eqref{eq:sigma3wave}--\eqref{eq:eps3wave} directly:
\begin{empheq}[box=\fbox]{align*}
\sig_x(t) &= \frac{E_b \Ab}{2 \As}\big[\eps_I(t) + \eps_R(t) + \eps_T(t)\big] \\
\epsdot_x(t) &= \frac{\Cz}{\Ls}\big[\eps_I(t) - \eps_R(t) - \eps_T(t)\big] \\
\eps_x(t) &= \int_0^t \epsdot_x(\tau)\,d\tau
\end{empheq}
This produces the dynamic stress--strain curve $\sig_x(\eps_x)$, which is the
primary deliverable of the test.

\subsection{Step 5: Lateral (Y, Z) Analysis}

Because no active pulse is applied to the Y/Z directions, those bars only record
the specimen's outward push from Poisson contraction and damage dilation. The
Y and Z output bars each show a single outgoing wave $\eps_Y(t)$ and $\eps_Z(t)$.
These give the dynamic component of lateral stress (added on top of the static
confining pre-stress $\sig_2$, $\sig_3$):
\begin{equation}
  \sig_y^{\text{dyn}}(t) = \frac{E_b \Ab}{\As}\,\eps_Y(t),
  \qquad
  \sig_z^{\text{dyn}}(t) = \frac{E_b \Ab}{\As}\,\eps_Z(t).
\end{equation}
Note: the factor of $1/2$ disappears compared to equation~\eqref{eq:sigma3wave}
because there is only one bar carrying force on each lateral face, not two
(the opposite face is held by the hydraulic cylinder, not by another bar).

\subsection{Validity Limits for Mode 1}

\begin{itemize}[leftmargin=*]
  \item Striker velocity $V \leq \SI{50}{m/s}$ (mechanical limit). At $V$ near
        \SI{50}{m/s}, $\sigpeak \approx \SI{1020}{MPa}$ already exceeds
        $\sigprop$ -- the striker cannot be fired at full velocity without
        bar yielding. Practical maximum is $V \approx \SI{40}{m/s}$.
  \item Static pre-stress range (current digital twin):
        $0 \leq \sig_1, \sig_2, \sig_3 \leq \SI{50}{MPa}$ per axis.
  \item Strain rate range: $\epsdot \approx \SIrange{50}{500}{s^{-1}}$.
\end{itemize}

% =====================================================
% MODE 2 -- EM UNIAXIAL
% =====================================================
\newpage
\section{Mode 2: Electromagnetic Uniaxial Impact}

This is the simplest mode of the new electromagnetic facility. The gas gun and
striker are replaced by an electromagnetic pulse generator, but only one
generator (along $+X$) is used. Apart from the pulse shape, the analysis is
identical to Mode~1.

\subsection{Hardware Difference From Mode 1}

\begin{itemize}[leftmargin=*]
  \item \textbf{Pulse generator}: capacitor bank discharged through a primary
        coil produces a Lorentz force on a flyer plate (or directly on the bar
        end), launching a stress pulse.
  \item \textbf{Bars}: square cross-section $\SI{50}{mm} \times \SI{50}{mm}$
        by default ($\Ab = \SI{2500}{mm^2}$). The square bars provide flat
        contact with cubic specimens. Bar geometry is editable in the Step~1
        sidebar (square or round, with user diameter).
  \item \textbf{Pre-stress}: servo-controlled hydraulic cylinders apply
        $\sig_1$, $\sig_2$, $\sig_3$ independently. The current digital twin
        uses the same \SI{50}{MPa} slider range as Mode~1; raise the slider
        cap in code if your physical facility supports a larger envelope.
\end{itemize}

\subsection{The Half-Sine Pulse Shape}

Where the gas-gun pulse is approximately rectangular, the electromagnetic
generator produces a clean half-sine pulse:
\begin{empheq}[box=\fbox]{equation}
  \sig_I(t) = \sigpeak \,\sin\!\left(\frac{\pi t}{\tau}\right),
  \qquad 0 \leq t \leq \tau.
  \label{eq:halfsine}
\end{empheq}
The peak amplitude $\sigpeak$ and duration $\tau$ are \emph{independently}
tuneable through capacitor charging voltage and circuit inductance, respectively.
This decoupling is the key advantage of the electromagnetic system: students can
sweep through a range of strain rates by varying $\tau$ alone, holding $\sigpeak$
fixed (or vice versa) -- a controlled experiment that is impossible with the
gas gun, where $\sigpeak$ and $\tau$ are linked through the striker.

\subsection{Why Half-Sine Pulses Are Better}

\begin{conceptbox}[Frequency content matters]
A perfectly rectangular pulse contains arbitrarily high-frequency Fourier
components, which travel at different speeds in the bar (Pochhammer--Chree
dispersion). This makes the pulse develop spurious oscillations as it
propagates. A half-sine pulse has most of its energy near the fundamental
frequency $f \sim 1/(2\tau)$, with rapidly decaying higher harmonics. Less
high-frequency content means less dispersion, cleaner signals, and more reliable
analysis.

A half-sine pulse also reaches dynamic stress equilibrium in the specimen faster
than a rectangular pulse of the same amplitude, because its smooth rise gives
the specimen time to reverberate internally before peak stress is reached. This
is critical for testing brittle materials, which would otherwise fail before
equilibrium and produce unreliable strength measurements.
\end{conceptbox}

\subsection{Step-by-Step Wave Analysis (Same as Mode 1)}

Because there is still only one driving bar, the equations are identical to
Mode~1. The student-facing procedure is:
\begin{stepbox}[EM uniaxial procedure]
\begin{enumerate}[leftmargin=*, itemsep=3pt]
  \item Apply static pre-stress $\sig_1$ (axial), $\sig_2$, $\sig_3$ (confining)
        via the three independent servo cylinders. Record baseline strain on all
        gauges. Wait for the system to settle (\SIrange{30}{60}{s}).
  \item Charge the capacitor bank to the voltage required for the desired
        $\sigpeak$. Set the inductance/circuit parameters for the desired $\tau$.
  \item Trigger the pulse generator. Capture all bar gauge signals at sampling
        rates $\geq \SI{1}{MHz}$.
  \item Subtract the pre-trigger baseline from each channel. (This is what
        removes the static pre-stress from the recorded signals -- see
        Section~\ref{subsec:prestress}.)
  \item Time-shift each channel to the specimen-face reference (incident gauge
        forward by $L_g/\Cz$, reflected backward by $L_g/\Cz$, etc.).
  \item Verify dynamic stress equilibrium: $|F_1 - F_2|/\bar F < 0.05$.
  \item Apply the three-wave equations~\eqref{eq:sigma3wave}--\eqref{eq:eps3wave}
        to obtain $\sig_x^{\text{dyn}}(t)$, $\epsdot_x(t)$, $\eps_x(t)$.
  \item Add back the static $\sig_1$ baseline to get total axial stress
        $\sig_x^{\text{total}}(t) = \sig_1 + \sig_x^{\text{dyn}}(t)$.
  \item Plot $\sig_x^{\text{total}}$ versus $\eps_x$. Read peak strength and
        average strain rate.
  \item Check that maximum bar stress remained below $\sigprop = \SI{930}{MPa}$.
\end{enumerate}
\end{stepbox}

\subsection{Worked Example: Sandstone with Triaxial Pre-stress}

A student wants to test sandstone (UCS $\approx \SI{80}{MPa}$, Mohr-Coulomb slope
$k_{\text{conf}} \approx 4.5$) under
$\sig_1 = \SI{20}{MPa}$ axial baseline and $\sig_2 = \sig_3 = \SI{50}{MPa}$
confining, at strain rate $\sim \SI{200}{s^{-1}}$. Estimate whether this is safe.

\textbf{Step 1 -- Predicted dynamic peak strength of the specimen.} Only the
\emph{lateral} pre-stresses $\sig_2, \sig_3$ enter the rock strength formula:
\begin{align*}
  \sigpeak^{\text{rock,dyn}}
    &= (\sig_{c0} + k_{\text{conf}} \cdot \max(\sig_2, \sig_3)) \cdot \text{DIF} \\
    &= (80 + 4.5 \times 50) \,\text{MPa} \cdot
       \big(1 + 0.18 \log_{10}(200)\big) \\
    &= 305 \cdot 1.41 \approx \SI{430}{MPa}.
\end{align*}

\textbf{Step 2 -- Total axial stress at peak.} Add the static axial baseline:
\begin{equation*}
  \sigpeak^{\text{total}} = \sig_1 + \sigpeak^{\text{rock,dyn}}
                          = 20 + 430 \approx \SI{450}{MPa}.
\end{equation*}

\textbf{Step 3 -- Required bar stress.} The bar carries only the dynamic
component (the static $\sig_1$ is held by the cylinder). Since $\As/\Ab = 1$
for a \SI{50}{mm} specimen between \SI{50}{mm} square bars,
\begin{equation*}
  \sig_{\text{bar}} = \sigpeak^{\text{rock,dyn}} \cdot \frac{\As}{\Ab}
                    \approx \SI{430}{MPa}.
\end{equation*}

\textbf{Conclusion:} \SI{430}{MPa} is well below the \SI{930}{MPa} bar limit.
Safe to proceed. Charge the capacitor for $\sigpeak \approx \SI{500}{MPa}$
(margin above 430 to drive the specimen through peak), choose
$\tau = \SI{200}{\micro s}$ to give a strain rate of order \SI{200}{s^{-1}}.
The simulator will report the total axial stress reaching $\sim$\SI{450}{MPa}
on the $\sig_X$ trace, which is the static $\sig_1 = \SI{20}{MPa}$ plus the
dynamic peak.

% =====================================================
% MODE 3 -- ASYNCHRONOUS TRIAXIAL
% =====================================================
\newpage
\section{Mode 3: Asynchronous Triaxial Impact}

This mode uses three of the six electromagnetic generators to apply pulses along
$+X$, $+Y$, and $+Z$, each with an independent time delay. It enables the study
of how the \emph{order} of multi-axis loading affects material failure -- a
question that classical Hopkinson testing cannot address.

\subsection{Hardware Configuration}

Three pulse generators fire pulses on three orthogonal incident bars. The other
three bars (at the $-X$, $-Y$, $-Z$ ends) act as transmission/output bars and
remain passive. In this mode, all three positive-axis directions carry their
own three-wave problem; the negative-axis directions only show transmitted
signals.

\subsection{The Time-Delay Concept}

Let pulse $i \in \{x, y, z\}$ be defined as:
\begin{equation}
  \sig_{I,i}(t) = \sigpeak \, \sin\!\left(\frac{\pi (t - \Delta t_i)}{\tau}\right),
  \qquad \Delta t_i \leq t \leq \Delta t_i + \tau.
  \label{eq:asyncpulse}
\end{equation}
Setting $\Delta t_x = 0$, $\Delta t_y = \SI{50}{\micro s}$,
$\Delta t_z = \SI{100}{\micro s}$ produces a sequential triaxial loading: the
specimen first experiences uniaxial dynamic compression along $X$, then
biaxial after \SI{50}{\micro s}, then triaxial after \SI{100}{\micro s}.

\begin{conceptbox}[Why ordering matters]
Brittle failure in rock and concrete is path-dependent. Cracks initiated in one
loading direction can be either suppressed (if confinement arrives quickly) or
allowed to propagate (if confinement arrives late). Asynchronous loading lets
us \emph{isolate} this effect from the synchronous case, producing data on the
fundamental question: do dynamic constitutive models need a stress-path term?

Steps~2--3 of the integrated workspace turn this informal question into a
quantitative one through the normalised-delay parameter
$\Delta t^\ast = \Delta t / t_{\text{travel}}$ and the regime-classifier map
described in Section~\ref{sec:wave_damage}.
\end{conceptbox}

\subsection{Three-Wave Analysis on Each Axis}

Because each axis carries its own incident wave, each axis is analysed with the
full three-wave method, separately:
\begin{empheq}[box=\fbox]{equation}
  \sig_i(t) = \frac{E_b \Ab}{2 \As}
    \big[\eps_{I,i}(t) + \eps_{R,i}(t) + \eps_{T,i}(t)\big],
  \quad i \in \{x, y, z\}.
\end{empheq}
\begin{empheq}[box=\fbox]{equation}
  \epsdot_i(t) = \frac{\Cz}{\Ls}
    \big[\eps_{I,i}(t) - \eps_{R,i}(t) - \eps_{T,i}(t)\big],
  \quad i \in \{x, y, z\}.
\end{empheq}
The full triaxial stress state of the specimen at time $t$ is then:
\begin{equation}
  \boldsymbol{\sig}(t) =
  \begin{bmatrix}
    \sig_1^0 + \sig_x^{\text{dyn}}(t) & 0 & 0 \\
    0 & \sig_2^0 + \sig_y^{\text{dyn}}(t) & 0 \\
    0 & 0 & \sig_3^0 + \sig_z^{\text{dyn}}(t)
  \end{bmatrix},
\end{equation}
where the superscript $0$ denotes the static pre-stress applied by the hydraulic
cylinders before the dynamic event.

\subsection{Step-by-Step Procedure}

\begin{stepbox}[Asynchronous triaxial procedure]
\begin{enumerate}[leftmargin=*, itemsep=3pt]
  \item Apply static pre-stress $\sig_1^0$ (axial), $\sig_2^0$, $\sig_3^0$
        (confining) via the three independent servo-cylinders. Record baselines.
  \item Choose desired pulse delays $\Delta t_y$, $\Delta t_z$ (relative to
        $\Delta t_x = 0$). Typical values:
        $\Delta t = 0$ (synchronous, ground truth), $\Delta t = \tau/4$
        (mild asynchrony), $\Delta t = \tau/2$ (strong asynchrony).
  \item Charge all three capacitor banks to identical voltages (for equal
        $\sigpeak$ on each axis).
  \item Trigger generators with the programmed delays. Synchronisation jitter
        must be $\ll \tau$; aim for $< \SI{1}{\micro s}$ on a \SI{200}{\micro s}
        pulse.
  \item Capture six bar signals (3 incident, 3 transmission) plus the four
        lateral output gauges (Y$_1$, Y$_2$, Z$_1$, Z$_2$) for equilibrium
        verification.
  \item Subtract baselines, time-align each channel to its own specimen face.
  \item Apply the three-wave method on each axis independently.
  \item Plot the principal stress trajectory $(\sig_1(t), \sig_2(t), \sig_3(t))$
        in 3D principal-stress space (or use the $p$--$q$ projection in Step~3).
  \item Compare with a matched synchronous test ($\Delta t = 0$) to extract the
        path-dependence effect.
\end{enumerate}
\end{stepbox}

\subsection{Validity Notes}

\begin{itemize}[leftmargin=*]
  \item All three axes are subject to the same bar plasticity limit. For
        $\sigpeak = \SI{500}{MPa}$ pulses on all three axes, no bar exceeds
        \SI{500}{MPa} -- safe.
  \item Equilibrium verification is harder in this mode because the specimen
        is responding simultaneously to three different loadings. Check
        $|F_{i,1} - F_{i,2}| / \bar F_i < 0.05$ on \emph{each} axis separately.
  \item The constitutive interpretation requires a fully 3D stress-path
        analysis; do not attempt to fit a uniaxial constitutive model to data
        from this mode.
\end{itemize}

% =====================================================
% MODE 4 -- SYMMETRIC SUPERPOSITION
% =====================================================
\newpage
\section{Mode 4: Symmetric Multidirectional Impact (Wave Superposition)}

This is the most powerful and most physically distinctive mode. All six
electromagnetic generators fire \emph{simultaneously}, with identical pulses on
opposing bar pairs $(+X, -X)$, $(+Y, -Y)$, $(+Z, -Z)$. The two pulses on each
axis arrive at the specimen at the same instant, superposing constructively in
force while cancelling in momentum. The specimen is \emph{squeezed} from all
sides simultaneously.

\begin{newfeature}[Why this mode is special]
\begin{itemize}[leftmargin=*, itemsep=2pt]
  \item \textbf{Force doubling at each face} via constructive superposition.
  \item \textbf{Zero net momentum transfer} -- the specimen does not accelerate.
  \item \textbf{No transmission bar exists} on the active axes; both
        $\pm$-direction bars are incident-only.
  \item \textbf{Bar stress is unchanged} from uniaxial mode at the same
        $\sigpeak$. The doubling occurs in the \emph{specimen} stress, not the
        bar stress -- a profound advantage for high-stress testing.
  \item \textbf{Pure hydrostatic loading} possible when all six fire equally:
        deviatoric stress $q \to 0$, the specimen probes its volumetric (cap)
        response.
\end{itemize}
\end{newfeature}

\subsection{The Geometry of Superposition}

Consider a single axis, say $X$. In this mode, both the $+X$ and $-X$ bars carry
identical incident pulses arriving at the specimen at the same instant. Define:
\begin{align}
  \eps_{I}^{+}(t) &= \text{incident wave on the } +X \text{ bar (positive-going)} \\
  \eps_{I}^{-}(t) &= \text{incident wave on the } -X \text{ bar (negative-going)}
\end{align}
By construction, in symmetric mode, $\eps_I^+ = \eps_I^- = \eps_I$ at the
specimen position. After the pulses interact with the specimen, each bar carries
a reflected wave $\eps_R^+$ or $\eps_R^-$ travelling back outward. The simulator
emits both $\eps_R^+$ and $\eps_R^-$ explicitly so that the symmetric energy
folding (Section~\ref{subsec:expdata}) can use them directly.

\subsection{Step-by-Step Derivation of the Symmetric Wave Equations}

We now derive the analogue of the three-wave equations for symmetric mode. The
derivation has four small steps; students should follow each one.

\subsubsection{Step A: Force at the specimen face from each bar}

The bar is still a 1D elastic medium. The force it exerts on the specimen face
is (using equation~\eqref{eq:F1} for the $+X$ side and the equivalent for $-X$):
\begin{align}
  F^{+}(t) &= E_b \Ab \big[\eps_I^{+}(t) + \eps_R^{+}(t)\big],
  \label{eq:Fplus}\\
  F^{-}(t) &= E_b \Ab \big[\eps_I^{-}(t) + \eps_R^{-}(t)\big].
  \label{eq:Fminus}
\end{align}
Each force pushes \emph{inward} on the specimen.

\subsubsection{Step B: Force balance on the specimen}

The specimen, treated as a rigid body to leading order, must satisfy
$F^{+} = F^{-}$ (otherwise it would accelerate). By symmetry of the input
pulses, this is automatic. The specimen stress is the average force divided by
specimen cross-section:
\begin{align}
  \sig_s(t) &= \frac{F^{+} + F^{-}}{2 \As}
            = \frac{E_b \Ab}{2 \As}\big[\eps_I^{+} + \eps_R^{+} + \eps_I^{-} + \eps_R^{-}\big].
\end{align}
Under perfect symmetry, $\eps_I^+ = \eps_I^- = \eps_I$ and
$\eps_R^+ = \eps_R^- = \eps_R$, so:
\begin{empheq}[box=\fbox]{equation}
  \sig_s^{\text{sym}}(t)
    = \frac{2 E_b \Ab}{2 \As}\big[\eps_I(t) + \eps_R(t)\big]
    = \frac{E_b \Ab}{\As}\big[\eps_I(t) + \eps_R(t)\big].
  \label{eq:sigsym}
\end{empheq}
Compare with the standard three-wave result~\eqref{eq:sigma3wave}: the prefactor
has \emph{doubled} (no factor of $1/2$) because we now have two bars driving
the specimen, not one.

\subsubsection{Step C: Velocity of the specimen face}

Consider the $+X$ specimen face. The $+X$ bar end velocity is, from~\eqref{eq:v1}:
\[
  v^{+}_{\text{bar}} = \Cz \big[\eps_I^{+} - \eps_R^{+}\big] = \Cz \big[\eps_I - \eps_R\big].
\]
This is the velocity at which the $+X$ face moves \emph{inward} (compressing the
specimen). By symmetry, the $-X$ face moves inward at the same speed. The total
\emph{rate of compression} of the specimen is the sum of these two inward
velocities divided by the specimen length:
\begin{empheq}[box=\fbox]{equation}
  \epsdot_s^{\text{sym}}(t)
    = \frac{v^{+} + v^{-}}{\Ls}
    = \frac{2\Cz}{\Ls}\big[\eps_I(t) - \eps_R(t)\big].
  \label{eq:epsdotsym}
\end{empheq}
The strain rate has \emph{doubled} compared to uniaxial mode. The factor of 2
in equation~\eqref{eq:epsdotsym} is geometric: both faces move inward, not just
one. The simulator further uses the simplified driver
$\epsdot_s^{\text{sym}}(t) \approx (2\Cz/\Ls)\,\eps_I^+(t)$ in its time-step
update to avoid a numerical deadlock at $t = 0$ where the reflected wave equals
the incident wave before any specimen stress has built up; this approximation
becomes exact once equilibrium is reached.

\subsubsection{Step D: Reflected wave from the specimen response}

To close the analysis, we need one more relation -- between $\eps_R$ and the
specimen response. From force balance at the $+X$ face,
$F^{+} = \sig_s \As \cdot \tfrac{1}{2}$ (each face carries half the specimen
force balance, since both faces drive equally). Substituting~\eqref{eq:Fplus}:
\[
  E_b \Ab \big[\eps_I + \eps_R\big] = \frac{1}{2} \sig_s \As.
\]
Solving for $\eps_R$:
\begin{equation}
  \eps_R(t) = \frac{\sig_s(t) \As}{2 E_b \Ab} - \eps_I(t).
  \label{eq:Rsym}
\end{equation}
This is the analogue of $\eps_T = \sig_s \As / (E_b \Ab)$ in the standard
analysis, but with a factor of 2 in the denominator because two bars share the
work.

\subsection{Summary Table: Mode 1 vs Mode 4 Equations}

\begin{center}
\renewcommand{\arraystretch}{1.4}
\begin{tabular}{@{}lcc@{}}
\toprule
\textbf{Quantity} & \textbf{Modes 1--3 (one-sided)} & \textbf{Mode 4 (symmetric)} \\
\midrule
Specimen stress & $\sig = \dfrac{E_b\Ab}{2\As}(\eps_I + \eps_R + \eps_T)$
                & $\sig = \dfrac{E_b\Ab}{\As}(\eps_I + \eps_R)$ \\
Strain rate     & $\epsdot = \dfrac{\Cz}{\Ls}(\eps_I - \eps_R - \eps_T)$
                & $\epsdot = \dfrac{2\Cz}{\Ls}(\eps_I - \eps_R)$ \\
Bars active     & $\eps_I, \eps_R, \eps_T$ (three bars) & $\eps_I, \eps_R$ (each of two bars) \\
Effective drive & $\sig_{\text{drive}} = \sigpeak$ & $\sig_{\text{drive}} = 2\sigpeak$ \\
Bar stress max  & $\sig_{\text{bar}}^{\max} \leq \sigpeak$
                & $\sig_{\text{bar}}^{\max} \leq \sigpeak$ \\
Specimen $\sig_{\max}$ achievable & $\sim \sigpeak$ & $\sim 2\sigpeak$ \\
\bottomrule
\end{tabular}
\end{center}

\begin{warningbox}[The fundamental advantage of symmetric mode]
The bar stress is the \emph{same} in symmetric mode as in uniaxial mode at the
same $\sigpeak$, but the specimen experiences \emph{double} the stress. This
means symmetric mode can reach specimen stresses that would require bar yielding
in any one-sided mode. For testing strong rocks under high confinement, this is
the only way to access the high-stress regime without damaging the bars.
\end{warningbox}

\subsection{The Three Symmetric Sub-Modes}

The simulator offers three configurations:

\begin{description}[leftmargin=2em, style=nextline]
  \item[Symmetric $\pm X$ only]
    Two pulses on $\pm X$; Y and Z bars passive. The specimen experiences
    enhanced uniaxial loading: doubled stress, doubled strain rate, but no
    triaxial state. Useful as a baseline against which the multi-axis modes
    are compared.

  \item[Symmetric $\pm X \pm Y$ (biaxial)]
    Four pulses on $\pm X$ and $\pm Y$; Z passive (carries only static
    confinement and Poisson reaction). The specimen experiences true biaxial
    dynamic squeeze. The deviatoric stress along Z (the axis of asymmetry)
    drives any shear-related failure.

  \item[Full hydrostatic ($\pm X, \pm Y, \pm Z$)]
    All six pulses fire identically. Under perfect symmetry,
    $\sig_x = \sig_y = \sig_z$, so the deviatoric stress
    $q = 0$ and the specimen cannot fail by shear yielding. Instead, it
    deforms volumetrically along its compaction cap. This probes the
    pressure-volume response, captured by the curve $p(\eps_v)$ where
    $p = (\sig_x + \sig_y + \sig_z)/3$ and $\eps_v = \eps_x + \eps_y + \eps_z$.
\end{description}

\subsection{Pressure--Volume Response in the Hydrostatic Sub-Mode}

For full hydrostatic loading, the relevant constitutive curve is the
\emph{compaction curve}, governed by the bulk modulus $K$ and an eventual
pressure cap $p_{\text{cap}}$ at which pore collapse and grain rearrangement
take over:
\begin{equation}
  \frac{dp}{d\eps_v} = K \quad \text{(initial elastic regime)},
  \qquad
  p \to p_{\text{cap}} \quad \text{(post-cap plateau)}.
\end{equation}
For typical porous rocks, $p_{\text{cap}} \sim 4 \cdot \text{UCS}$ at moderate
strain rate. The simulator switches to this constitutive regime automatically
when full hydrostatic mode is selected.

\subsection{Step-by-Step Symmetric-Mode Procedure}

\begin{stepbox}[Symmetric multidirectional procedure]
\begin{enumerate}[leftmargin=*, itemsep=3pt]
  \item Apply static pre-stress: set $\sig_1^0 = \sig_2^0 = \sig_3^0$ for the
        hydrostatic sub-mode, or set independent values for the $\pm X$ or
        $\pm X \pm Y$ sub-modes.
  \item Choose $\sigpeak$ (single-bar peak amplitude). Remember the specimen
        will see $2 \sigpeak$ on each active axis. Verify the bar limit:
        $\sigpeak < \sigprop = \SI{930}{MPa}$.
  \item Choose $\tau$ for the desired strain rate. In symmetric mode, the
        strain rate is doubled compared to uniaxial mode at the same $\tau$,
        so for a target $\epsdot$, set $\tau \approx 2 \cdot \tau_{\text{uniaxial}}$.
  \item Charge all six (or four, or two) capacitor banks to the same voltage
        for identical $\sigpeak$ on each active bar.
  \item Synchronise the trigger pulses. Synchronisation jitter must be
        $\ll \tau/10$ for clean superposition; modern trigger electronics
        achieve $< \SI{0.5}{\micro s}$.
  \item Capture all bar gauge signals (4 per active axis: incident-wave gauge
        and reflected-wave gauge on each of the two opposing bars).
  \item Verify symmetry: $\eps_I^+ \approx \eps_I^-$ and
        $\eps_R^+ \approx \eps_R^-$ on each axis. Asymmetry indicates either
        synchronisation failure or non-uniform specimen damage; in either
        case, the analysis is invalid.
  \item Apply the symmetric-mode equations~\eqref{eq:sigsym}
        and~\eqref{eq:epsdotsym} to obtain $\sig_i(t)$ and $\epsdot_i(t)$ on
        each active axis.
  \item For hydrostatic sub-mode: compute mean pressure $p(t)$ and volumetric
        strain $\eps_v(t)$, plot $p(\eps_v)$.
  \item For biaxial sub-mode: extract both the active-axis $\sig$--$\eps$
        curves and the passive-axis Poisson response.
  \item Confirm bar plasticity envelope is satisfied:
        $\max_i |\sig_{\text{bar},i}(t)| < \sigprop$.
\end{enumerate}
\end{stepbox}

\subsection{A Worked Symmetric Example}

A student wants to test concrete (UCS \SI{40}{MPa}) under fully hydrostatic
dynamic loading with target peak pressure \SI{800}{MPa}. Static pre-stress
$\sig_1^0 = \sig_2^0 = \sig_3^0 = \SI{50}{MPa}$ (the slider maximum).

\textbf{Required single-pulse amplitude:} since the specimen sees
$2 \sigpeak$ on each axis, and the loading is hydrostatic so the mean pressure
equals the axial stress:
\begin{align*}
  p_{\text{peak}} &= \sig_1^0 + 2\sigpeak \\
  \SI{800}{MPa} &= \SI{50}{MPa} + 2\sigpeak \\
  \sigpeak &= \SI{375}{MPa}.
\end{align*}
\textbf{Bar stress check:} each bar carries one pulse at $\sigpeak$:
\[
  \sig_{\text{bar}}^{\max} \approx \SI{375}{MPa} < \SI{930}{MPa} \;\checkmark
\]
Plenty of margin -- this configuration is safely within the bar limit, even
though it would require $\sigpeak = \SI{750}{MPa}$ (very near the bar limit) in
uniaxial mode.

\textbf{Pulse duration:} for target $\epsdot \sim \SI{300}{s^{-1}}$ in
symmetric mode, the doubled strain-rate factor means we want
$\tau \approx 2\Cz \eps_I / (\Ls \cdot \epsdot_{\text{target}}) \sim$
\SI{200}{\micro s}.

\textbf{Result:} Charge each of six capacitor banks for $\sigpeak = \SI{375}{MPa}$,
$\tau = \SI{200}{\micro s}$. Fire synchronously with $\sig_{\text{conf}} =
\SI{50}{MPa}$ pre-applied. Expected outcome: pure compaction response, no shear
failure, bar limits respected.

% =====================================================
% PART 4 -- COMPARISON & GUIDANCE
% =====================================================
\newpage
\section{Comparing the Four Modes Side by Side}

\begin{center}
\renewcommand{\arraystretch}{1.5}
\small
\begin{tabular}{@{}p{3.6cm}p{2.6cm}p{2.6cm}p{2.6cm}p{2.6cm}@{}}
\toprule
\textbf{Property} & \textbf{Mode 1: Gas-Gun} & \textbf{Mode 2: EM Uniaxial} & \textbf{Mode 3: EM Async} & \textbf{Mode 4: Symmetric} \\
\midrule
Pulse shape & Trapezoidal (shaped) & Half-sine & Half-sine $\times 3$ & Half-sine $\times 6$ \\
Active bars & 1 (X only) & 1 (X only) & 3 (+X, +Y, +Z) & Up to 6 ($\pm$XYZ) \\
$\sig_1, \sig_2, \sig_3$ slider range & 0--50 MPa each & 0--50 MPa each & 0--50 MPa each & 0--50 MPa each \\
Pulse independence & $\sigpeak \propto V$, $\tau \propto L$ & Independent & Independent & Independent \\
Specimen $\sig_{\max}$ & $\sim \sigpeak$ & $\sim \sigpeak$ & $\sim \sigpeak$ each axis & $\sim 2\sigpeak$ each axis \\
Bar stress at limit & $\sigpeak$ & $\sigpeak$ & $\sigpeak$ & $\sigpeak$ \\
Stress path control & Uniaxial only & Uniaxial only & Yes (asynchronous) & Hydrostatic / biaxial \\
Pure hydrostatic? & No & No & Approximately if $\Delta t = 0$ & Yes \\
Best for & UCS, basic SHPB & Clean rate-dependence & Path-dependent failure & Cap mechanics, high stress \\
\bottomrule
\end{tabular}
\end{center}

\noindent\textit{Note on slider range:} the \SI{50}{MPa} pre-stress cap reflects
the current digital twin's calibration. The slider maximum is set in one place
(\texttt{max\_conf} in the simulator file) and can be raised if a physical bench
upgrade supports a larger envelope.

\subsection{Choosing the Right Mode for a Research Question}

\begin{itemize}[leftmargin=*, itemsep=3pt]
  \item \textbf{Mode 1 (Gas-Gun)} when: the research is a baseline uniaxial
        test, or a comparison against published SHPB literature. Reliable,
        well-understood, but limited in confinement and pulse shape.

  \item \textbf{Mode 2 (EM Uniaxial)} when: rate-dependent constitutive
        behaviour is the primary question. Use the independent control of
        $\sigpeak$ and $\tau$ to sweep $\epsdot$ at fixed $\sig$ and vice
        versa.

  \item \textbf{Mode 3 (EM Asynchronous)} when: the question is whether
        loading order matters. Particularly relevant for blast-induced
        rockburst and seismic loading, where the multi-axis pulses arrive
        sequentially. Pair this mode with the Step~2 regime classifier and
        the Step~3 $p$--$q$--$\theta$ path.

  \item \textbf{Mode 4 (EM Symmetric)} when: (a) the target specimen stress
        exceeds the bar plasticity limit in one-sided modes, or (b) the
        research targets the volumetric (cap) response of the material, which
        is inaccessible in any one-sided mode.
\end{itemize}

\subsection{Common Pitfalls and How the Simulator Helps}

\begin{itemize}[leftmargin=*]
  \item \textbf{Bar plasticity violation.} The simulator's three-tier warning
        system (caution / warning / critical) catches this in real time. Plan
        every test in the simulator first.
  \item \textbf{Loss of stress equilibrium.} For very brittle specimens or
        very short pulses, the specimen may fail before reverberations
        equilibrate it. Check $|F_1 - F_2| / \bar F < 0.05$. The Step~2
        ``Wave model'' tab reports the equilibrium window
        $t_{\text{eq}} \in [3, 5]\,t_{\text{travel}}$ as a coloured band.
  \item \textbf{Confusing strain rate with strain.} A short, high-amplitude
        pulse delivers high $\epsdot$ but possibly insufficient total strain
        to fail the specimen. Check both metrics in the summary panel.
  \item \textbf{Asynchronous mode misinterpretation.} Do not fit a uniaxial
        constitutive model to data from Mode 3. Use a 3D stress-path-aware
        model (e.g., Drucker--Prager with cap, or Mohr--Coulomb with
        non-associated flow) -- Step~3 provides the $p$--$q$--$\theta$
        diagnostics needed to identify the path.
  \item \textbf{Symmetric-mode synchronisation failure.} If
        $\eps_I^+ \neq \eps_I^-$ on the recorded gauges, the assumed symmetry
        is broken and the doubled-stress analysis no longer applies.
  \item \textbf{Sign-convention mismatch on uploaded gauges.} The Step~1
        reducer assumes textbook signs ($\eps_I > 0$, $\eps_R < 0$,
        $\eps_T > 0$ for a compression test). Use the \emph{Compression
        convention} radio button to flip data acquired with the opposite
        polarity; an inverted convention typically appears as a negative
        absorbed energy, which the simulator now flags explicitly.
\end{itemize}

% =====================================================
% INTEGRATED WORKSPACE
% =====================================================
\newpage
\section{The Integrated Workspace}

The Virtual Tri-HB is now delivered as a single integrated Streamlit
application, \texttt{tri\_hb\_integrated.py}, that bundles four analysis
steps under one shared experimental setup. This section walks through how to
launch the application and how to navigate the four-step workflow. The two
sections that follow (Sections~\ref{sec:expdata} and~\ref{sec:wave_damage})
document the specific equations used by the experimental reducer and by the
wave/damage workspace.

\subsection{Launching the Application}

\begin{stepbox}[Local launch (development / classroom laptop)]
\begin{enumerate}[leftmargin=*, itemsep=3pt]
  \item Install Python $\geq 3.9$.
  \item Open a terminal in the folder containing \texttt{tri\_hb\_integrated.py},
        \texttt{Tri-HB.py}, \texttt{wave\_damage.py}, and
        \texttt{requirements.txt}.
  \item Install dependencies: \texttt{pip install -r requirements.txt}.
  \item Launch the integrated app: \texttt{streamlit run tri\_hb\_integrated.py}.
  \item The browser opens at \texttt{http://localhost:8501}. If it does not,
        open that URL manually.
\end{enumerate}
\end{stepbox}

For shared use across a class, push the repository to GitHub and connect it at
\texttt{share.streamlit.io}. Students then access the integrated workspace
through a public URL with no local installation.

\subsection{Top-Level Navigation: Five Pages}

A single radio in the sidebar selects which page is active. The intent is to
move through them in order, but jumping back is supported because each step
reads the shared \texttt{tri\_hb\_latest\_config} and
\texttt{tri\_hb\_latest\_result} from session state.

\begin{description}[leftmargin=2em, style=nextline]
  \item[Overview]
    Landing page that explains the workflow. Use it to confirm the manual
    matches the deployed app version.

  \item[Step 1: Setup, simulator and data]
    Defines the shared experimental setup (material, specimen, bar, loading)
    and produces either a \emph{simulated} stress--strain history or a
    \emph{reduced} stress--strain history from uploaded gauge data. A second
    radio inside the sidebar switches between two tasks: \emph{Test design
    and simulator} (the four-mode simulator, Sections~2--5) and
    \emph{Experimental data analysis} (the bar-gauge / direct stress-strain
    reducer, Section~\ref{sec:expdata}).

  \item[Step 2: Wave model]
    Pulse timing, travel time, equilibrium window, and the wave-superposition
    / damage-memory regime classifier. Loads the
    \texttt{wave\_damage.py} module with the ``Wave model'' tab first.

  \item[Step 3: Stress path and analysis]
    $p$--$q$--$\theta$ paths, invariant histories, comparison against
    the reduced experimental data from Step~1. Same module, with the ``Stress
    path'' tab first.

  \item[Step 4: Damage model and DEM validation]
    Damage evolution, stiffness and wave-speed degradation, energy balance,
    DEM/experimental descriptors, parametric study, and export. Same module,
    with the ``Damage evolution'' tab first.
\end{description}

\subsection{Step 1 Sidebar: Two Tabs}

Inside the Step~1 simulator task, the sidebar is split into two tabs that
match the operational separation a real test team uses.

\begin{description}[leftmargin=2em, style=nextline]
  \item[Experimental setup]
    Specimen material (preset + editable $E$, UCS, $\nu$, $\rho$), specimen
    geometry (cross-section, side / diameter, loading length $\Ls$), and bar
    properties (bar $E_b$, $\Cz$, cross-section). Defaults match the
    42CrMo square 50$\times$50\,mm bench; every value can be overridden
    to model a different facility.

  \item[Test loading]
    Loading mode (4 modes), symmetric-axis selector (if Mode~4), pulse settings
    (striker velocity for gas-gun; $\sigpeak$ and $\tau$ for EM), static
    pre-stress sliders for $\sig_1$, $\sig_2$, $\sig_3$, and the asynchronous
    timing sliders if Mode~3 is active.
\end{description}

When Step~1 is configured for ``Experimental data analysis'' instead of the
simulator, the sidebar replaces the test-loading controls with the data-source
selector and the column-mapping controls described in
Section~\ref{sec:expdata}.

\subsection{Main Panel Tabs in the Step 1 Simulator}

The main panel updates live as inputs change. Five tabs are available; the
fourth ($p$--$\eps_v$) appears only in symmetric mode.

\begin{description}[leftmargin=2em, style=nextline]
  \item[Bar Waveforms] Time-history of all strain gauge signals: incident,
    reflected, transmitted, and lateral output bars. In symmetric mode, both
    $\eps_I^+$ and $\eps_I^-$ are plotted to verify synchronisation.

  \item[Specimen Stresses] Total $\sig_X$, $\sig_Y$, $\sig_Z$ on the specimen.
    In symmetric mode, the mean pressure $p(t)$ is overlaid as a dashed
    reference. A horizontal dashed line at \SI{930}{MPa} shows the bar
    proportional limit for visual safety check.

  \item[$\sig$--$\eps$ Curve] Loading branch only. The curve traces from the
    start of the pulse through peak strength and post-peak softening, then
    stops when strain plateaus.

  \item[$p$--$\eps_v$] Visible only in symmetric mode. The compaction
    (volumetric) curve shows initial bulk-modulus slope followed by a cap
    pressure plateau.

  \item[Equations] Reference card showing the wave-analysis equations for
    the active mode, including the standard three-wave method and the
    symmetric-mode variant.
\end{description}

\subsection{Bar Plasticity Warning Banner}

Whenever the maximum bar stress approaches the proportional limit, a coloured
warning banner appears at the top of the page:

\begin{itemize}[leftmargin=*, itemsep=2pt]
  \item \textcolor{warnorange}{\textbf{Caution}} (yellow tint) at
        $0.85\,\sigprop$.
  \item \textcolor{warnorange}{\textbf{Warning}} (orange) at $\sigprop$.
        Wave analysis becomes nonlinear; results are unreliable.
  \item \textcolor{dangerred}{\textbf{Critical}} (red) at the yield stress
        \SI{1080}{MPa}. Bars deform plastically; the test is invalid and
        should not be run on the physical apparatus.
\end{itemize}

This is the simulator's primary safety feature. \textbf{Always check the banner
is clear (or at most cautionary) before considering the corresponding
configuration for physical testing.}

\subsection{Data Export from Step 1}

Four export buttons sit at the bottom of the simulator page. Filenames are
auto-generated from the configuration so re-running with different parameters
produces uniquely-named files for batch comparison.

\begin{center}
\renewcommand{\arraystretch}{1.4}
\small
\begin{tabular}{@{}p{3.6cm}p{2cm}p{8.5cm}@{}}
\toprule
\textbf{Button} & \textbf{Format} & \textbf{Contents} \\
\midrule
All Signals & CSV &
  Time-series: all bar gauge strains
  ($\eps_I^{\pm}, \eps_R^{\pm}, \eps_T, \eps_Y, \eps_Z$ as available),
  specimen stresses ($\sig_X, \sig_Y, \sig_Z$), strains, pressure,
  volumetric strain, axial strain rate. \\
$\sig$--$\eps$ Curve & CSV &
  Loading branch only: axial strain, axial stress, mean pressure,
  volumetric strain. \\
Configuration & JSON &
  Full input config + summary metrics + bar/rock properties + warning
  state + ISO 8601 export timestamp. Provides full provenance. \\
Full workbook & XLSX &
  Multi-sheet Excel: Sheet~1 Signals, Sheet~2 Stress-Strain,
  Sheet~3 Configuration. \\
\bottomrule
\end{tabular}
\end{center}

\begin{conceptbox}[Recommended export workflow]
For coursework and research projects:
\begin{enumerate}[leftmargin=2em, itemsep=2pt]
  \item Run the simulator with your chosen configuration.
  \item Verify the bar plasticity banner is clear.
  \item Export the JSON \emph{and} either the CSV or XLSX, both filed under
        the auto-generated filename.
  \item Move to Step~2--4 to interpret wave timing, stress path, and damage,
        which automatically inherit the Step~1 setup through session state.
  \item Export Step~4's CSV for the damage and energy histories alongside the
        Step~1 file. The JSON provides the audit trail linking both.
\end{enumerate}
\end{conceptbox}

\subsection{A Worked Streamlit Session}

\begin{stepbox}[Session: Granite at high confinement, EM Symmetric XYZ]
\begin{enumerate}[leftmargin=*, itemsep=3pt]
  \item In the top sidebar, select \emph{Step 1: Setup, simulator and data}
        and leave the inner radio on \emph{Test design and simulator}.
  \item Open the \emph{Experimental setup} tab. Choose \emph{Material preset =
        granite}; keep the auto-filled $E$, UCS, $\nu$. Set \emph{Specimen
        cross-section = Square/cube}, \emph{Side length = \SI{50}{mm}},
        \emph{Loading length = \SI{50}{mm}}. Leave bar settings on the
        \SI{50}{mm} square default.
  \item Switch to the \emph{Test loading} tab. Choose \emph{Loading mode = EM
        Symmetric Multi-axis}; the axis selector appears -- pick \emph{Full
        hydrostatic ($\pm X \pm Y \pm Z$)}.
  \item Set \emph{Single-pulse peak amplitude = \SI{400}{MPa}}, \emph{Pulse
        duration = \SI{200}{\micro s}}. Set
        $\sig_1 = \sig_2 = \sig_3 = \SI{50}{MPa}$.
  \item Inspect the warning banner. If it shows ``Caution'' or above, reduce
        $\sigpeak$ until the banner is clear.
  \item Click the \emph{$p$--$\eps_v$} tab. Read the slope at low $\eps_v$
        (the dynamic bulk modulus $K_d$) and the plateau at high $\eps_v$
        (the cap pressure).
  \item Click \emph{Full workbook (XLSX)}; verify the \emph{Configuration}
        sheet records all your inputs.
  \item Move to Step~2: confirm the equilibrium band and check
        $\Delta t^\ast$ (this should be 0 for a synchronous test).
  \item Move to Step~3: check the $p$--$q$ path collapses onto the
        $q \approx 0$ axis (hydrostatic).
  \item Move to Step~4: read the final damage $D$, the central damage
        fraction $D_c$, and the energy descriptors. Export the CSV for
        downstream comparison with DEM or experimental data.
\end{enumerate}
\end{stepbox}

% =====================================================
% EXPERIMENTAL DATA ANALYSIS
% =====================================================
\newpage
\section{Step 1 (Alternate Task): Experimental Data Analysis}
\label{sec:expdata}
\label{subsec:expdata}

When the inner Step~1 radio is set to ``Experimental data analysis'', the
simulator's loading controls are replaced by a reducer that converts raw bar
gauge or direct stress-strain data into the same set of outputs the simulator
produces. This means students can compare a real measurement and a simulation
on the same axes, using the same downstream pages (Steps~2--4).

\subsection{Three Data Sources}

\begin{description}[leftmargin=2em, style=nextline]
  \item[Latest simulator result]
    Reuses the most recent simulator output stored in session state. Useful
    for verifying that the reducer and the simulator agree, and for handing
    simulated data into the energy balance, wave model, and damage page.

  \item[Upload bar strain gauges]
    A CSV/XLSX with columns for time and the incident, reflected, and
    transmitted bar strains. The student maps each column manually and picks
    the unit (\si{\micro\strain} or strain), the time unit
    (\si{\micro s}, ms, or s), and the compression convention (positive or
    negative). The reducer then applies the SHPB equations and produces
    stress, strain, strain rate, and energy histories.

  \item[Upload direct stress-strain]
    A CSV/XLSX that already contains a reduced stress and strain. Only unit
    conversion and a derivative for the strain rate are needed; this branch
    is intended for comparison against the simulator output.
\end{description}

\subsection{Equations Used by the Reducer}

Let $E_b$ be the bar Young's modulus, $\Cz = \sqrt{E_b/\rho_b}$ the bar wave
speed, $\Ab$ the bar cross-section, $\As$ and $\Ls$ the specimen area and
length, and $\eps_I$, $\eps_R$, $\eps_T$ the incident, reflected, and
transmitted bar strains.

\paragraph{Specimen stress (three-wave average).}
\begin{equation}
  \sig_s(t) = \frac{E_b \Ab}{2 \As}\,\bigl[\eps_I(t)+\eps_R(t)+\eps_T(t)\bigr].
\end{equation}

\paragraph{Specimen stress (transmitted-wave only).}  Assumes force equilibrium
$\eps_I + \eps_R \approx \eps_T$:
\begin{equation}
  \sig_s(t) = \frac{E_b \Ab}{\As}\,\eps_T(t).
\end{equation}

\paragraph{Specimen strain rate and strain.}
\begin{align}
  \epsdot_s(t) &= \frac{\Cz}{\Ls}\,\bigl[\eps_I(t)-\eps_R(t)-\eps_T(t)\bigr], \\
  \eps_s(t)   &= \int_0^t \epsdot_s(\tau)\,d\tau \quad \text{(trapezoidal)}.
\end{align}
Under perfect equilibrium $\eps_I + \eps_R = \eps_T$ the rate reduces to
$\epsdot_s = -2\Cz\,\eps_R/\Ls$; this app keeps the general 3-wave form so
pre-equilibrium signals are still handled.

\paragraph{Bar-wave energies.}
\begin{align}
  W_I(t) &= E_b \Ab \Cz \int_0^t \eps_I^2\,d\tau, \\
  W_R(t) &= E_b \Ab \Cz \int_0^t \eps_R^2\,d\tau, \\
  W_T(t) &= E_b \Ab \Cz \int_0^t \eps_T^2\,d\tau, \\
  W_{\text{abs}}(t) &= W_I(t) - W_R(t) - W_T(t).
\end{align}
The units check is
$\Ab\,[\si{m^2}] \cdot E_b\,[\si{Pa}] \cdot \Cz\,[\si{m/s}]
\cdot \int \eps^2\,dt\,[\si{s}] = \si{J}$.

\paragraph{Symmetric-mode (Mode~4) energy folding.}
Opposing bar pairs contribute both incident pulses and both reflected pulses
explicitly:
\begin{equation}
  W_I = E_b \Ab \Cz \int \bigl(\eps_{I,+}^{\,2} + \eps_{I,-}^{\,2}\bigr)\,dt,
  \qquad
  W_R = E_b \Ab \Cz \int \bigl(\eps_{R,+}^{\,2} + \eps_{R,-}^{\,2}\bigr)\,dt.
\end{equation}

\paragraph{Absorbed energy from simulator stress-strain.}
\begin{equation}
  W_{\text{abs}}(t) = \int_0^t \sig_s(\tau)\,\epsdot_s(\tau)\,V_s\,d\tau,
  \qquad V_s = \As\,\Ls.
\end{equation}
The simulator stores $\sig_s$ as the \emph{total} axial stress (static pre-load
+ dynamic), so the reflected, transmitted, and absorbed energies for a
simulator-derived history are closed against the available incident energy.
This monotone closure prevents the static-preload contribution from appearing
as spurious wave energy and keeps the cumulative budget consistent with the
incident envelope.

\paragraph{Direct stress-strain branch.}
Only unit conversion and a derivative:
\begin{equation}
  \epsdot_s(t) = \frac{d\eps_s}{dt} \quad \text{(numpy central difference)}.
\end{equation}
Strain unit options: strain (1), percent (\(\div 100\)), microstrain
(\(\times 10^{-6}\)). Stress unit options: MPa (passthrough), Pa
(\(\div 10^6\)). The compression-convention selector applies a single global
sign so the formulas above remain consistent with the textbook convention
$\eps_I > 0$, $\eps_R < 0$, $\eps_T > 0$ for a compression test.

\subsection{Reducer Outputs and Energy Sanity Check}

The reducer writes the result to session state as
\texttt{tri\_hb\_reduced\_data} so that Steps~2--4 can use it for validation
overlays. Outputs include time, stress, strain, strain rate, and (for bar-gauge
data) the four energy histories.

A built-in energy sanity check raises a warning whenever any of:
\[
  W_T > W_I, \qquad W_R > W_I, \qquad W_{\text{abs}} = W_I - W_R - W_T < 0
\]
hold for uploaded bar-gauge data. The most common cause is a sign or unit
mismatch in the uploaded file (e.g.\ microstrain reported as strain, or
compression positive in some columns and negative in others). The warning
reports the numerical excess so the student knows how serious the mismatch is.

\subsection{A Worked Reducer Session}

\begin{stepbox}[Session: Comparing a measurement against the simulator]
\begin{enumerate}[leftmargin=*, itemsep=3pt]
  \item In Step~1, configure and run the simulator to match the geometry and
        loading of the physical test. Note the simulator's peak stress and
        strain rate.
  \item Switch the inner Step~1 radio to ``Experimental data analysis''.
  \item In the sidebar, choose \emph{Data source = Upload bar strain gauges}.
        The bar / specimen geometry inputs inherit the simulator's setup
        automatically; override them only if the upload was taken with a
        different rig.
  \item Upload the CSV/XLSX. Map the time column, $\eps_I$, $\eps_R$,
        $\eps_T$. Pick the time and strain units; check the compression
        convention against the raw data (textbook: $\eps_I > 0$ for
        compression).
  \item Choose \emph{Stress reduction = Three-wave average} for a baseline,
        then switch to \emph{Transmitted-wave only} to confirm the two
        formulas agree once equilibrium is reached.
  \item Read the four metric cards: peak stress, peak strain, peak strain
        rate, absorbed energy. Inspect the time-history and energy tabs; if
        the energy warning appears, fix the sign or unit before continuing.
  \item Move to Step~3. The ``DEM/experimental descriptors'' overlay (in
        Step~4) will use this reduced dataset for validation comparisons.
\end{enumerate}
\end{stepbox}

% =====================================================
% WAVE, STRESS PATH, DAMAGE, DEM PAGE
% =====================================================
\newpage
\section{Steps 2--4: Wave Model, Stress Path, Damage and DEM Validation}
\label{sec:wave_damage}

Steps~2, 3, and 4 are three views into the same underlying module
(\texttt{wave\_damage.py}). They share the experimental setup defined in
Step~1 (material, specimen length, prestress, pulse amplitude, pulse duration,
loading dimensionality, asynchronous delay) and switch only which tab is
opened first. The equations below are listed once and apply across all three
steps.

\subsection{Wave Model Equations (Step 2 Tab)}

\paragraph{Elastic constants and wave speeds.}
\begin{align}
  M &= \frac{E(1-\nu)}{(1+\nu)(1-2\nu)}, \\
  c_p &= \sqrt{M/\rho}, \qquad c_s = \sqrt{G/\rho},
\end{align}
with $M$ the constrained (P-wave) modulus. Wave travel time and the equilibrium
window are
\begin{equation}
  t_{\text{travel}} = \Ls/c_p, \qquad t_{\text{eq}} \in [3,\,5]\,t_{\text{travel}}.
\end{equation}
The equilibrium window is shaded on the time-history plot.

\paragraph{Normalised delay and regime classifier.}
\begin{equation}
  \Delta t^\ast = \Delta t / t_{\text{travel}}.
\end{equation}
The simulator partitions $\Delta t^\ast$ into four regimes:
\begin{center}
\renewcommand{\arraystretch}{1.25}
\begin{tabular}{@{}lll@{}}
\toprule
$\Delta t^\ast$ & Regime & Dominant physics \\
\midrule
$< 1$ & Synchronous wave superposition & Constructive overlap at specimen \\
$1$ to $3$ & Reverberation-coupled & Internal echoes still strong \\
$3$ to $10$ & Transitional & Decaying reverberations \\
$> 10$ & Sequential, damage-memory & Each pulse acts on damaged state \\
\bottomrule
\end{tabular}
\end{center}

The Step~2 transition map also overlays two dimensionless guide curves,
\begin{align}
  C_{\text{wave}}(\Delta t^\ast) &= \exp\!\left[-(\Delta t^\ast/1.2)^2\right], \\
  C_{\text{damage}}(\Delta t^\ast) &=
    \frac{1}{1+\exp\!\left[-(\Delta t^\ast-5.0)/1.3\right]}.
\end{align}
These curves are heuristic dominance indicators, not calibrated damage laws. The
third legend item in the app is the vertical marker for the current
$\Delta t^\ast$ value, so the map contains two curves plus one current-case
marker.

\paragraph{Superposition factor at the specimen centre.}
\begin{equation}
  \sig_{\text{centre}}(t) = \sig_{x0} + A_x\,g\!\left(t - \tfrac{\Ls}{2c_p}\right)
   + A_x r\,g\!\left(t - \Delta t - \tfrac{\Ls}{2c_p}\right),
\end{equation}
\begin{equation}
  \eta_{\text{sup}} = \frac{\max_t \sig_{\text{centre}}(t) - \sig_{x0}}{A_x},
\end{equation}
where $g(\cdot)$ is the chosen pulse envelope (Hann, half-sine, or rectangular)
and $r$ is the right-bar amplitude ratio. A value $\eta_{\text{sup}} \to 2$
indicates full constructive overlap; $\eta_{\text{sup}} \to 1$ indicates
sequential pulses.

\subsection{Stress-Path Equations (Step 3 Tab)}

Assuming the principal axes align with the bar axes, the diagonal Cauchy
tensor gives the invariants
\begin{align}
  p &= \tfrac{1}{3}(\sig_x + \sig_y + \sig_z), \\
  q &= \sqrt{\tfrac{1}{2}\bigl[(\sig_x - \sig_y)^2 + (\sig_y - \sig_z)^2 + (\sig_z - \sig_x)^2\bigr]}, \\
  J_2 &= \tfrac{1}{2}(s_1^2 + s_2^2 + s_3^2), \quad s_i = \sig_i - p, \\
  J_3 &= s_1 s_2 s_3, \\
  \cos(3\theta) &= \frac{3\sqrt{3}}{2}\,\frac{J_3}{J_2^{3/2}}.
\end{align}
The Lode angle $\theta$ distinguishes triaxial compression
($\theta = 0$) from triaxial extension ($\theta = 60^\circ$) and from pure
shear ($\theta = 30^\circ$).

\paragraph{Failure surface and failure index.}
\begin{align}
  h(\theta) &= 1 + a_\theta\bigl(1 - \cos 3\theta\bigr), \\
  q_f(p,\theta) &= \bigl(A + B\,p^{n}\bigr)\,h(\theta), \\
  F(t) &= q(t) / q_f(t).
\end{align}
The failure index $F$ is unitless; $F > 1$ means the stress state lies outside
the static failure surface and damage is allowed to grow.

\paragraph{Equivalent strain rate} (elastic estimate from diagonal stresses):
\begin{equation}
  \eps_i = \sig_i / E, \quad \epsdot_i = d\eps_i / dt,
\end{equation}
\begin{equation}
  \epsdot_{\text{eq}} = \sqrt{\tfrac{2}{3} \cdot \tfrac{1}{2}\bigl[(\epsdot_x - \epsdot_y)^2 + (\epsdot_y - \epsdot_z)^2 + (\epsdot_z - \epsdot_x)^2\bigr]}.
\end{equation}
This is a planning-grade indicator built from elastic strains only -- it
ignores plastic strain and lateral coupling. Use it as a regime tag, not as a
constitutive strain rate.

\subsection{Damage and Energy Equations (Step 4 Tab)}

\paragraph{Damage evolution} (Perzyna-type, saturating, rate-sensitive):
\begin{equation}
  \dot D = \frac{(1-D)^{\alpha}}{\tau_D}
           \Bigl\langle \frac{F-1}{F_0} \Bigr\rangle^{m}
           \Bigl(\frac{|\epsdot_{\text{eq}}|}{\epsdot_0}\Bigr)^{\beta},
  \qquad D \in [0,1],
\end{equation}
where $\langle\cdot\rangle$ denotes Macaulay brackets (damage grows only when
$F > 1$).

\paragraph{Stiffness and wave-speed degradation.}
\begin{equation}
  E(D) = E_0 (1 - D), \qquad c_p(D) = c_{p0}\,\sqrt{\max(1 - D,\,0)}.
\end{equation}

\paragraph{Energy indicators} (per unit volume):
\begin{align}
  W_{\text{el}} &= \frac{1}{2E}\bigl[\sig_x^2 + \sig_y^2 + \sig_z^2
                   - 2\nu(\sig_x\sig_y + \sig_y\sig_z + \sig_z\sig_x)\bigr], \\
  W_{\text{input}}(t) &= \int_0^t \bigl(\sig_x \epsdot_x + \sig_y \epsdot_y + \sig_z \epsdot_z\bigr)\,d\tau, \\
  W_{\text{diss}} &\approx W_{\text{input}} - W_{\text{el}} + W_{\text{el}}(0).
\end{align}
With stresses in MPa and dimensionless strains, $W$ is in MJ/m\textsuperscript{3}.

\paragraph{Synthetic DEM/experimental descriptors.}
The damage profile $D(x)$ across the specimen is constructed from three
contributions (left wave, right wave, and central superposition), then summary
statistics are derived:
\begin{align}
  D_c &= \frac{\int_{\text{centre}} D(x)\,dx}{\int_0^{\Ls} D(x)\,dx}
       \quad \text{(central damage fraction)}, \\
  S_x &= 1 - \frac{|D_{\text{left}} - D_{\text{right}}|}{D_{\text{left}} + D_{\text{right}}}
       \quad \text{(symmetry index)}.
\end{align}
These are model indicators for planning and interpretation, not actual DEM
outputs -- compare against a real DEM simulation or CT scan to validate.

\subsection{Sidebar Controls Used Across Steps 2--4}

\begin{description}[leftmargin=2em, style=nextline]
  \item[Control level] \emph{Guided} or \emph{Advanced}. Advanced exposes the
    failure parameters $A, B, n$, the Lode amplitude $a_\theta$, the damage
    time scale $\tau_D$, the saturation exponent $\alpha$, the overstress
    exponent $m$, the rate exponent $\beta$, the reference strain rate
    $\epsdot_0$, and the normalisation $F_0$.

  \item[Use Step 1 setup] When ticked (default if Step~1 has been run), the
    material, geometry, prestress, pulse amplitude, pulse duration, and
    loading dimensionality are read from the shared session state. Untick it
    to enter a standalone setup.

  \item[Analysis goal] \emph{Quick validation}, \emph{Stress-path focus}, or
    \emph{Damage calibration}. Pre-fills sensible defaults for the chosen
    workflow.
\end{description}

\subsection{Export from Step 4}

A single CSV download bundles every time-series produced by the wave/damage
module: input pulses, $\sig_x$ / $\sig_y$ / $\sig_z$, the invariants $p$, $q$,
$\theta$, the failure surface $q_f$, the failure index $F$, the equivalent
strain rate, $D$, $\dot D$, $E(D)$, $c_p(D)$, $W_{\text{input}}$,
$W_{\text{el}}$, and $W_{\text{diss}}$.

% =====================================================
% APPENDIX
% =====================================================
\newpage
\appendix
\section{Reference Tables}

\subsection{Bar properties (42CrMo steel) -- default values}

These are the simulator's default bar parameters; each is editable in the
Step~1 \emph{Experimental setup} sidebar tab.

\begin{center}
\begin{tabular}{@{}lcc@{}}
\toprule
\textbf{Quantity} & \textbf{Symbol} & \textbf{Default value} \\
\midrule
Young's modulus & $E_b$ & \SI{210}{GPa} \\
Density & $\rho_b$ & \SI{7850}{kg/m^3} \\
Bar wave speed & $\Cz$ & \SI{5172}{m/s} \\
Proportional limit & $\sigprop$ & \SI{930}{MPa} \\
Yield stress (approx.) & $\sigma_Y$ & \SI{1080}{MPa} \\
Round bar area (gas-gun, $\Phi 40$) & $\Ab$ & \SI{1257}{mm^2} \\
Square bar area ($\SI{50}{mm} \times \SI{50}{mm}$) & $\Ab$ & \SI{2500}{mm^2} \\
\bottomrule
\end{tabular}
\end{center}

\subsection{Default rock parameters in the simulator}
\begin{center}
\renewcommand{\arraystretch}{1.3}
\begin{tabular}{@{}lcccccc@{}}
\toprule
\textbf{Rock} & $E_s$ (GPa) & $\sig_{c0}$ (MPa) & $\nu$ & $b_{\text{rate}}$ & $k_{\text{conf}}$ & $\eps_{\text{peak}}$ \\
\midrule
Sandstone & 15 & 80 & 0.20 & 0.18 & 4.5 & 0.008 \\
Granite   & 50 & 180 & 0.25 & 0.12 & 5.5 & 0.005 \\
Concrete  & 30 & 40 & 0.20 & 0.22 & 4.0 & 0.006 \\
\bottomrule
\end{tabular}
\end{center}

Selecting a preset in the Step~1 sidebar populates $E_s$, $\sig_{c0}$, and
$\nu$ from this table; the values are then editable so non-listed materials
can be modelled.

\subsection{Default damage-model parameters (Step 4 Advanced mode)}
\begin{center}
\renewcommand{\arraystretch}{1.3}
\begin{tabular}{@{}lcc@{}}
\toprule
\textbf{Symbol} & \textbf{Meaning} & \textbf{Default} \\
\midrule
$A$ & Failure-surface constant & \SI{15}{MPa} \\
$B$ & Failure-surface pressure coefficient & 1.3 \\
$n$ & Failure-surface pressure exponent & 0.75 \\
$a_\theta$ & Lode-angle amplitude & 0.10 \\
$\tau_D$ & Damage time scale & \SI{30}{\micro s} \\
$\alpha$ & Saturation exponent & 1.0 \\
$m$ & Overstress exponent & 2.0 \\
$\beta$ & Rate exponent & 0.15 \\
$\epsdot_0$ & Reference strain rate & \SI{1.0}{s^{-1}} \\
$F_0$ & Failure-index normalisation & 1.0 \\
\bottomrule
\end{tabular}
\end{center}

\subsection{Useful identities}
\begin{align*}
  \text{Acoustic impedance:}\quad & Z = \rho \Cz \\
  \text{Pulse-to-bar transmission:}\quad
    & \frac{F_{\text{transmitted}}}{F_{\text{incident}}}
      = \frac{2 Z_2}{Z_1 + Z_2} \\
  \text{Dynamic Increase Factor (rate hardening):}\quad
    & \text{DIF} = 1 + b_{\text{rate}} \log_{10}\!\frac{\epsdot}{\epsdot_{\text{ref}}} \\
  \text{Mean pressure:}\quad & p = (\sig_1 + \sig_2 + \sig_3)/3 \\
  \text{Deviatoric stress:}\quad
    & q = \sqrt{\tfrac{1}{2}\big[(\sig_1 - \sig_2)^2 + (\sig_2 - \sig_3)^2 + (\sig_3 - \sig_1)^2\big]} \\
  \text{Lode angle:}\quad
    & \cos(3\theta) = (3\sqrt{3}/2)\,J_3/J_2^{3/2}
\end{align*}

\section{Suggested Student Exercises}

\begin{enumerate}[leftmargin=*, itemsep=4pt]
  \item In Mode 1 (Gas-Gun), find the maximum striker velocity for which a
        granite specimen with \SI{50}{MPa} confinement remains within the bar
        proportional limit. Verify your hand calculation against the
        simulator's warning system.
  \item In Mode 2 (EM Uniaxial), keep $\sigpeak = \SI{400}{MPa}$ fixed and
        vary $\tau$ from \SI{50}{\micro s} to \SI{300}{\micro s}. Plot peak specimen
        stress vs strain rate. Extract the rate hardening parameter
        $b_{\text{rate}}$ and compare with the rock's nominal value.
  \item In Mode 3 (Async), run sandstone at $\sig_{\text{conf}} = \SI{30}{MPa}$
        with $\Delta t_y = \Delta t_z = 0$ (synchronous), then with
        $\Delta t_y = \SI{100}{\micro s}$, $\Delta t_z = \SI{200}{\micro s}$.
        Compare peak $\sig_x$ between the two cases, then open Step~2 and
        record $\Delta t^\ast$ and the regime classifier output. Discuss the
        link between $\Delta t^\ast$ and the change in peak strength.
  \item In Mode 4 (Symmetric XYZ), find the peak specimen pressure achievable
        with single-pulse $\sigpeak = \SI{900}{MPa}$ (right at the bar limit).
        Compare with the maximum pressure achievable in Mode 2 (EM Uniaxial)
        before bar yielding.
  \item Derive the symmetric-mode equations~\eqref{eq:sigsym}--\eqref{eq:epsdotsym}
        from first principles (force and velocity continuity), without
        consulting Section~5.2.
  \item Extend the symmetric-mode derivation to the case where the two
        opposing pulses have \emph{different} amplitudes, say $\sigpeak^+$
        and $\sigpeak^-$. Show that the specimen experiences both a squeeze
        (proportional to $\sigpeak^+ + \sigpeak^-$) and a net acceleration
        (proportional to $\sigpeak^+ - \sigpeak^-$). At what amplitude ratio
        does this degenerate into Mode 2 (one-sided) behaviour?
  \item In Mode 2 (EM Uniaxial), test sandstone with $\sigpeak = \SI{300}{MPa}$,
        $\tau = \SI{200}{\micro s}$, and $\sig_2 = \sig_3 = \SI{30}{MPa}$.
        Run the simulation three times with $\sig_1 = 0$, $\SI{25}{MPa}$,
        and $\SI{50}{MPa}$. Export the $\sig$--$\eps$ CSV from each run.
        Overlay the three curves in your preferred plotting software and
        verify that increasing $\sig_1$ shifts the entire curve upward by
        a constant offset \emph{without} changing peak strength. Explain
        why $\sig_1$ behaves differently from $\sig_2$ and $\sig_3$.
  \item In Step~1, run any Mode~2 case. Then switch the inner radio to
        ``Experimental data analysis'' and choose \emph{Data source = Latest
        simulator result}. Confirm that the three-wave reduction reproduces
        the simulator's stress--strain curve. Switch to \emph{Transmitted-wave
        only} and check that the two reductions agree once stress equilibrium
        is reached.
  \item Take the same Mode~2 simulator result and download the All Signals
        CSV. Re-upload it through the ``Upload bar strain gauges'' branch
        with column mapping. Verify the reduced curve matches the
        simulator-derived curve, then deliberately mis-set the strain unit
        (e.g. claim ``microstrain'' for data actually stored as strain) and
        observe how the energy sanity warning fires.
  \item Move to Step~3 with any asynchronous Mode~3 simulator result. Read
        $\Delta t^\ast$ and the regime classifier off the Wave model tab,
        then trace the $p$--$q$ path in the Stress path tab. Identify the
        instant at which $F(t) = q(t)/q_f(t)$ first exceeds unity and
        compare with the time at which damage starts to grow on the Damage
        evolution tab in Step~4.
  \item Choose any test case that interests you and use the export buttons
        to download all four file types (Signals CSV, $\sig$--$\eps$ CSV,
        Configuration JSON, Full XLSX). Open the JSON file in a text editor
        and identify: (a) the simulator version, (b) the export timestamp,
        (c) the bar plasticity warning state, and (d) the peak specimen
        stress. Discuss how this metadata supports reproducibility in a
        research context.
\end{enumerate}

\vspace{1cm}
\begin{center}
{\large\itshape\color{deepblue} -- End of Manual --}\\[6pt]
\small\color{midblue} For corrections or suggestions, please contact the Tri-HB
research group.
\end{center}

\end{document}
